% ============================================================
% Cheat sheet, argomento 13: quale diagramma usare (dispensa, capitoli 6-9)
% ============================================================
\argomento{Quale diagramma?}{Dispensa: capitoli 6--9. Dalla domanda dell'esercizio al diagramma giusto}[Quale
diagramma]

\begin{guida}{Che cosa chiede l'esercizio}
\centering
\renewcommand{\arraystretch}{1.35}
\begin{tabular}{@{}p{8.6cm} l l@{}}
  \textbf{La domanda} & \textbf{Il diagramma} & \textbf{Argomento} \\ \hline
  Chi usa il sistema e che cosa può farci? Quali sono i requisiti funzionali? & casi d'uso & \vv{attore} \\
  Quali concetti esistono nel dominio, con quali dati e legami? & classi (dominio) & \vv{classe} \\
  Come sono fatte le classi Java del progetto? Disegna il diagramma di questo codice. & classi (progetto) & \vv{campisemplici} \\
  Chi chiama chi, e in che ordine, in uno scenario? Che cosa fa questo metodo? & sequenza & \vv{partecipante} \\
  Come reagisce un oggetto agli eventi? Ha delle fasi? Che cosa fa un'interfaccia grafica? & stato & \vv{stato} \\
  Qual è il flusso di lavoro, chi fa che cosa, che cosa avviene in parallelo? & attività & \vv{azione} \\
  Come si divide il codice in moduli? Quali dipendenze ci sono fra i moduli? & package & \vv{package} \\
  Quali pezzi compongono il sistema e quali interfacce offrono? & componenti & \vv{componente} \\
  Su quali macchine gira il sistema, e come sono collegate? & deployment & \vv{nodo} \\
\end{tabular}
\end{guida}

\begin{guida}{I diagrammi di UML 2: due famiglie}
\begin{itemize}[nosep, leftmargin=*]
  \item \textbf{Statici} (strutturali): \textbf{com'è fatto} il sistema. Classi, package, componenti, deployment (e oggetti).
  \item \textbf{Dinamici} (comportamentali): \textbf{che cosa fa} nel tempo. Casi d'uso, stato, attività e quelli di
        \textbf{interazione}: sequenza (il più usato), comunicazione, interaction overview, timing.
\end{itemize}
UML 2 ha 13 tipi di diagrammi; i più usati sono quello delle classi e quello di sequenza. Vale sempre la regola delle
\textbf{informazioni soppresse}: ciò che un diagramma non mostra non è per forza assente.
\end{guida}

\begin{guida}{Quelli che si confondono}
\centering
\renewcommand{\arraystretch}{1.35}
\begin{tabular}{@{}>{\bfseries}p{3.2cm} p{14.6cm}@{}}
  Stato o attività? & Lo \textbf{stato} descrive \textbf{un oggetto} che aspetta eventi e cambia comportamento
    (rettangoli = situazioni in cui l'oggetto \textit{si trova}). L'\textbf{attività} descrive un \textbf{processo} che
    avanza da solo (rettangoli = azioni che \textit{si fanno}), anche fra più oggetti (swimlane). \\
  Sequenza o attività? & La \textbf{sequenza} mostra \textbf{chi parla con chi} in un solo scenario; l'\textbf{attività}
    mostra la \textbf{logica} (decisioni, cicli, parallelismo). Per un algoritmo, meglio l'attività. \\
  Dominio o progetto? & Nel modello di \textbf{dominio} (concettuale): niente metodi (al più responsabilità), niente
    tipi e visibilità, niente verso di navigazione, niente «Sistema» né «id». Nel \textbf{progetto}: tutto questo
    c'è, e il diagramma corrisponde al codice. \\
  Classi o oggetti? & Il diagramma delle \textbf{classi} vale per tutte le istanze (molteplicità); un diagramma di
    \textbf{oggetti} \oltre\ fotografa un momento, con i valori dei campi. \\
  Componenti o deployment? & I \textbf{componenti} dicono che cosa c'è e chi offre quale interfaccia;
    il \textbf{deployment} dice su quale macchina gira. \\
\end{tabular}
\end{guida}

\begin{attenzione}{Le regole che valgono in tutti i diagrammi}
\begin{itemize}[nosep, leftmargin=*]
  \item Nomi delle classi al \textbf{singolare}; attori sostantivi, casi d'uso e azioni con un \textbf{verbo}.
  \item Le frecce \textbf{tratteggiate} con la punta aperta sono dipendenze («include», «extend», «call», «create»,
        «use», le dipendenze fra package e fra componenti).
  \item Un diagramma serve a \textbf{comunicare}: se è troppo affollato va diviso o semplificato (prevalgono le
        regole descrittive).
  \item I diagrammi dello stesso sistema devono essere \textbf{coerenti}: ogni messaggio di un diagramma di sequenza
        è un'operazione nel diagramma delle classi, ogni evento di un diagramma di stato è un metodo.
\end{itemize}
\end{attenzione}
